% Refactor-twin of `claim_3_14_proof_AddingInterventionNodes.tex` for
% refactor `eqViaNodeMap_injective`.  The refactor strengthens the
% predicate `eqViaNodeMap` (from claim_3_7) to `refactor_eqViaNodeMap`
% by adding a 1st conjunct `Set.InjOn f (G.J ∪ G.V)` on the carrier
% map; the four image-equality conjuncts (J, V, E, L) and the
% mathematical content of the lemma are unchanged.  This twin's only
% additions versus the existing `cdmg_typed_edges` twin are:
%
%   (i) a new paragraph in the statement block clarifying that
%       sub-claim (a)'s componentwise equality is read \emph{modulo the
%       canonical flatten map $\flattenIntExt$} and that the
%       strengthened reading requires $\flattenIntExt$ to be injective
%       on the iterated-extension's joint input-output carrier (the
%       only statement-level change; the four CDMG-component readings
%       and the LN-side prose are unchanged); and
%
%   (ii) a new ``Injectivity of the canonical flatten map
%        $\flattenIntExt$ on the iterated extended graph's $J \cup V$''
%        paragraph inserted into each of sub-claims (a-1) and (a-2),
%        just before each closing ``Combining \ldots''/``Re-running
%        \ldots'' line.  Mathematically each new paragraph is a
%        four-cell case analysis on the iterated graph's joint
%        input-output node set, with the disjointness hypothesis
%        $W_1 \cap W_2 = \emptyset$ consumed exactly in the cell-3
%        vs.~cell-4 cross-collision step (and its symmetric
%        counterpart).  The new paragraphs live at the LN/Lean
%        interface and are not load-bearing for the LN-level
%        mathematics of the lemma; they underwrite the strengthened
%        Lean-encoding of the ``componentwise equality up to canonical
%        carrier bijection'' predicate `refactor_eqViaNodeMap`
%        introduced by the root refactor (claim_3_7).
%
% Sub-claim (b) is \emph{not} touched: both sides of (b) live on the
% same single-extension carrier $J_{\doit(I_{W_1})} \cup V_{\doit(I_{W_1})}$
% (since the hard-intervention $\doit(W_2)$ preserves its argument
% CDMG's carrier), so there is no carrier-bijection identification in
% play and the LN-level equality is read as literal componentwise
% equality.  In Lean this is captured by the separate
% `addInterventionNodes_comm_hardIntervention` theorem (literal `=` of
% `CDMG (IntExtNode Node)`), which this refactor does not touch.
%
% Phase-7 cleanup renames this twin file over the original.

\documentclass[main]{subfiles}

\providecommand{\flattenIntExt}{\mathrm{flattenIntExt}}

\begin{document}
% ref: claim_3_14 (proof, with statement restated for readability)
% title: AddingInterventionNodes
\def\rowref{claim\_3\_14}
\phantomsection\label{claim_3_14}

% --- statement (restated) ---------------------------------------------------
\begin{Lem}[Adding intervention nodes commutes with disjoint hard interventions]\label{adding-intervention-nodes-commutes-with-disjoint-hard-interventions}
    Let $G = (J, V, E, L)$ be a CDMG (def \ref{def-cdmg}); throughout we use the notation of def \ref{not-cdmg} and recall the typing and axioms of def \ref{def-cdmg}, namely $J \cap V = \emptyset$, $E \subseteq (J \cup V) \times V$, $L \subseteq V \times V$, $L$ irreflexive ($(v_1, v_2) \in L \Rightarrow v_1 \neq v_2$), and $L$ symmetric ($(v_1, v_2) \in L \Leftrightarrow (v_2, v_1) \in L$). Let
    \[ W_1 \subseteq J \cup V \quad\text{and}\quad W_2 \subseteq J \cup V \]
    be two subsets of nodes of $G$ (overlap with $J$ is permitted on either side, matching the preconditions of both def \ref{def:G_hard_intervention} and def \ref{def:cdmg_intervention_nodes}), subject to the disjointness side condition
    \[ W_1 \cap W_2 \;=\; \emptyset. \]
    All three quantifiers ``for every CDMG $G$'', ``for every $W_1 \subseteq J \cup V$'', and ``for every $W_2 \subseteq J \cup V$ with $W_1 \cap W_2 = \emptyset$'' are read \emph{universally}; in particular the (vacuously disjoint) case $W_1 = W_2 = \emptyset$ is admitted, and so are the corner cases $W_1 \subseteq J$ or $W_2 \subseteq J$ (in which the corresponding intervention-node addition $\doit(I_{W_i})$ reduces to the identity by def \ref{def:cdmg_intervention_nodes} since $W_i \sm J = \emptyset$, and the corresponding hard intervention $\doit(W_i)$ acts only via its input-side relabelling and the removal of any edges into $W_i \cap V = \emptyset$ by def \ref{def:G_hard_intervention}).

    \emph{Disambiguation of the overloaded $\doit(\cdot)$ notation.}
    The notation $G_{\doit(\cdot)}$ used in this lemma is overloaded; it is disambiguated by the \emph{kind} of its argument, as follows. For any subset $W \subseteq J \cup V$:
    \begin{itemize}
        \item $G_{\doit(W)}$ denotes the \emph{hard intervention} on $W$ (def \ref{def:G_hard_intervention}): the CDMG with components $J_{\doit(W)} := J \cup W$, $V_{\doit(W)} := V \sm W$, $E_{\doit(W)} := E \sm \lC (v_1, v_2) \in E \st v_2 \in W \rC$, and $L_{\doit(W)}$ obtained from $L$ by removing every bidirected edge incident to a node of $W$ (item iv.\ of def \ref{def:G_hard_intervention}, read with the symmetrisation needed for $L_{\doit(W)}$ to satisfy the symmetry axiom of def \ref{def-cdmg}). No new nodes are introduced.
        \item $G_{\doit(I_W)}$ denotes the \emph{extended CDMG} obtained by adding intervention nodes for $W$ (def \ref{def:cdmg_intervention_nodes}): the CDMG with components $J_{\doit(I_W)} := J \cup \lC I_w \st w \in W \sm J \rC$, $V_{\doit(I_W)} := V$, $E_{\doit(I_W)} := E \cup \lC (I_w, w) \st w \in W \sm J \rC$, and $L_{\doit(I_W)} := L$, where $\lC I_w \st w \in W \sm J \rC$ is the family of fresh intervention symbols introduced by def \ref{def:cdmg_intervention_nodes} (type-disjoint from $J \cup V$ and pairwise type-disjoint), and where the notational shorthand $I_j := j$ for $j \in J \cap W$ means no new node is introduced on the $J \cap W$ branch.
    \end{itemize}
    The two constructions are distinguished by the kind of the argument under $\doit$: a node set $W \subseteq J \cup V$ (hard intervention, no new nodes) versus the corresponding intervention-node symbol $I_W$ (extension, fresh intervention nodes added for $W \sm J$).

    \emph{Well-typedness of the iterated operations on both sides.}
    Each iterated operation appearing in the two displayed equations below is well-typed per its constructor's precondition:
    \begin{itemize}
        \item By items i.\ and ii.\ of def \ref{def:cdmg_intervention_nodes}, the carrier of $G_{\doit(I_{W_i})}$ is $J_{\doit(I_{W_i})} \cup V_{\doit(I_{W_i})} = (J \cup V) \cup \lC I_w \st w \in W_i \sm J \rC \supseteq J \cup V$ (canonical inclusion via the ``unsplit'' map $\iota \colon J \cup V \hookrightarrow J_{\doit(I_{W_i})} \cup V_{\doit(I_{W_i})}$, $v \mapsto v$, as in \ref{claim_3_13}). Hence any $W_j \subseteq J \cup V$ satisfies $W_j \subseteq J_{\doit(I_{W_i})} \cup V_{\doit(I_{W_i})}$ (via $\iota$), so both $\lp G_{\doit(I_{W_i})} \rp_{\doit(I_{W_j})}$ and $\lp G_{\doit(I_{W_i})} \rp_{\doit(W_j)}$ are defined (per def \ref{def:cdmg_intervention_nodes} and def \ref{def:G_hard_intervention} respectively).
        \item By items i.\ and ii.\ of def \ref{def:G_hard_intervention} together with $W_i \subseteq J \cup V$, the carrier of $G_{\doit(W_i)}$ is $J_{\doit(W_i)} \cup V_{\doit(W_i)} = (J \cup W_i) \cup (V \sm W_i) = J \cup V$ as a set. Hence any $W_j \subseteq J \cup V$ satisfies $W_j \subseteq J_{\doit(W_i)} \cup V_{\doit(W_i)}$, so both $\lp G_{\doit(W_i)} \rp_{\doit(W_j)}$ and $\lp G_{\doit(W_i)} \rp_{\doit(I_{W_j})}$ are defined.
        \item The right-hand-side $G_{\doit(I_{W_1 \cup W_2})}$ is defined per def \ref{def:cdmg_intervention_nodes}, since $W_1 \cup W_2 \subseteq J \cup V$.
    \end{itemize}

    \emph{Definition of the mixed-argument notation $G_{\doit(I_{W_1}, W_2)}$ (introduced in this lemma).}
    The expression $G_{\doit(I_{W_1}, W_2)}$ appearing in the second displayed equation below is introduced here for the first time; it is not given by either of the two upstream definitions (def \ref{def:G_hard_intervention} takes a single node set, def \ref{def:cdmg_intervention_nodes} takes a single node set's worth of intervention symbols). We define it explicitly as the composition of the intervention-node-addition operation $\doit(I_{W_1})$ (def \ref{def:cdmg_intervention_nodes}) and the hard-intervention operation $\doit(W_2)$ (def \ref{def:G_hard_intervention}):
    \[ G_{\doit(I_{W_1}, W_2)} \;:=\; \lp G_{\doit(I_{W_1})} \rp_{\doit(W_2)}. \]
    Sub-claim (b) below asserts that the alternative ordering $\lp G_{\doit(W_2)} \rp_{\doit(I_{W_1})}$ produces the same CDMG, so this definition is order-independent. The comma in ``$\doit(I_{W_1}, W_2)$'' is a separator between two semantically distinct operations (intervention-node addition for $W_1$ on the left, hard intervention on $W_2$ on the right) and does \emph{not} denote a set-union $\doit(I_{W_1} \cup W_2)$, an indexed-set $\doit(I_{W_1, W_2})$, or any other parsing of the comma.

    \emph{The conclusion.} The following two equalities of CDMGs hold.
    \begin{enumerate}[label=(\alph*)]
        \item \emph{Adding intervention nodes commutes with itself (disjoint $W_1, W_2$).}
              The three CDMGs $\lp G_{\doit(I_{W_1})} \rp_{\doit(I_{W_2})}$, $\lp G_{\doit(I_{W_2})} \rp_{\doit(I_{W_1})}$, and $G_{\doit(I_{W_1 \cup W_2})}$ coincide:
              \[ \lp G_{\doit(I_{W_1})} \rp_{\doit(I_{W_2})} \;=\; \lp G_{\doit(I_{W_2})} \rp_{\doit(I_{W_1})} \;=\; G_{\doit(I_{W_1 \cup W_2})}. \]
        \item \emph{Adding intervention nodes commutes with hard intervention (disjoint $W_1, W_2$).}
              The two CDMGs $\lp G_{\doit(I_{W_1})} \rp_{\doit(W_2)}$ and $\lp G_{\doit(W_2)} \rp_{\doit(I_{W_1})}$ coincide and equal the mixed-notation CDMG $G_{\doit(I_{W_1}, W_2)}$ just defined:
              \[ \lp G_{\doit(I_{W_1})} \rp_{\doit(W_2)} \;=\; \lp G_{\doit(W_2)} \rp_{\doit(I_{W_1})} \;=\; G_{\doit(I_{W_1}, W_2)}. \]
              (Equivalently, the equality $\lp G_{\doit(I_{W_1})} \rp_{\doit(W_2)} = \lp G_{\doit(W_2)} \rp_{\doit(I_{W_1})}$ is the genuine content; the third term is a recap of the defining equation of $G_{\doit(I_{W_1}, W_2)}$.)
    \end{enumerate}

    \emph{Carrier mismatch in sub-claim (a) and the strengthened componentwise reading.}
    The two sides of sub-claim (a) have formally distinct node-carrier sets: the LHS $\lp G_{\doit(I_{W_1})} \rp_{\doit(I_{W_2})}$ lives on the iterated-extension carrier (carrying two layers of fresh intervention symbols introduced by the inner and outer applications of $\doit(I_\cdot)$), while the RHS $G_{\doit(I_{W_1 \cup W_2})}$ lives on the single-extension carrier (a single layer of fresh symbols indexed by $(W_1 \cup W_2) \sm J$). The LN identifies the two carriers set-theoretically via the canonical \emph{flatten map} $\flattenIntExt$ (made explicit in the proof body below: it sends each twice-tagged iterated-extension node to its single-step counterpart, collapsing the LN's ``$\iota(\iota(v)) = v$'' and ``outer-$I$ of inner-$\iota$ collapses to a single-step $I$'' conventions), and the equality in (a) is read as \emph{componentwise equality after applying $\flattenIntExt$ to the LHS}. To prevent this image-level reading from silently degenerating to a many-to-one collapse (which would equate iterated-extension CDMGs whose carriers differ in node count from the single-step CDMG), we additionally require that $\flattenIntExt$ be \emph{injective on the iterated-extension's carrier set} $J_{\mathrm{iter}} \cup V_{\mathrm{iter}}$ (where $\mathrm{iter}$ stands for the iterated-extension graph appearing on the LHS of (a-1) or (a-2)). This injectivity is the substantive content of the carrier identification: without it, two formally distinct iterated-extension nodes could share the same single-step image, and the LN's reading of (a) as a CDMG identity (not merely an image-level agreement) would fail. Sub-claim (b), by contrast, has no carrier mismatch: both $\lp G_{\doit(I_{W_1})} \rp_{\doit(W_2)}$ and $\lp G_{\doit(W_2)} \rp_{\doit(I_{W_1})}$ live on the single-extension carrier $J_{\doit(I_{W_1})} \cup V_{\doit(I_{W_1})}$ (the outer hard intervention $\doit(W_2)$ preserves its argument CDMG's carrier per def \ref{def:G_hard_intervention}, and the inner $\doit(I_{W_1})$ produces the same single-extension carrier in both orderings), so the equality is read as \emph{literal componentwise equality} with no canonical-bijection identification.

    Each equality in (a) and (b) is understood componentwise on the four components $(J, V, E, L)$ of def \ref{def-cdmg}: equal input-node sets, equal output-node sets, equal directed-edge sets, and equal bidirected-edge sets; the conjunctive unpacking into these four field-by-field equalities is deferred to the proof. For sub-claim (a) the four equalities are read modulo the canonical flatten map $\flattenIntExt$ (with the load-bearing injectivity requirement above); for sub-claim (b) the four equalities are read literally.
\end{Lem}

% --- proof ------------------------------------------------------------------
% Proof body adapted from the LN's \Claude{...} block at
% lecture-notes/lecture_notes/graphs.tex lines 845-874 via the
% `cdmg_typed_edges` twin, expanded as follows:
% (i) the structure now tracks the rewritten statement's (a) / (b) split, with
%     (a) decomposed into sub-claims (a-1), (a-2) and the swap-by-transitivity
%     step made explicit;
% (ii) for both (a-1) and the two orderings in (b), all four CDMG components
%     (J, V, E, L) are verified explicitly -- the LN dismisses V and L as
%     "unchanged"; we spell this out so a reader does not have to backsolve;
% (iii) the symmetric statement (a-2) is derived via "by symmetry of the roles
%       of W_1 and W_2" followed by commutativity of union, the swap then
%       recovered by transitivity through the shared right-hand side;
% (iv) the identification of the two orderings in (b) with the mixed-notation
%      CDMG G_{\doit(I_{W_1}, W_2)} is performed by appeal to the defining
%      equation of that notation introduced in the statement block;
% (v) the disjointness hypothesis W_1 \cap W_2 = \emptyset is consumed
%     explicitly at each step that uses it, so the verifier can check that no
%     silent dependence on disjointness sneaks into the V or L components;
% (vi) a new "Injectivity of the canonical flatten map $\flattenIntExt$ on the
%      iterated extended graph's $J \cup V$" paragraph is inserted at the end
%      of each of sub-claims (a-1) and (a-2)'s componentwise verification, just
%      before the closing "Combining ..." / "Re-running ..." lines (this is the
%      new content versus the `cdmg_typed_edges` twin -- the carrier-map
%      injectivity certificate that underwrites the strengthened Lean predicate
%      `refactor_eqViaNodeMap` introduced by refactor
%      `eqViaNodeMap_injective`).  Mathematically each new paragraph is a
%      four-cell case analysis using the disjointness hypothesis
%      W_1 \cap W_2 = \emptyset to rule out the (3)-vs-(4) cross-collision
%      (and its symmetric counterpart for (a-2)); the new paragraphs are not
%      load-bearing for the LN-level mathematics of the lemma.
\begin{proof}
We establish sub-claims (a) and (b) in turn; each is an equality of CDMGs (def \ref{def-cdmg}), verified componentwise on each of the four components $(J, V, E, L)$. For sub-claim (a) we additionally verify the carrier-map injectivity of the canonical flatten map $\flattenIntExt$ on the iterated-extension's joint input-output node set, as required by the strengthened componentwise reading recorded in the statement block above (carrier-mismatch paragraph). The well-typedness of every iterated operation appearing below is licensed by the corresponding bullet of the ``well-typedness'' paragraph in the statement block, which we cite without restating. Throughout, we freely use the explicit four-component spell-outs of the single-step CDMGs $G_{\doit(I_W)}$ (def \ref{def:cdmg_intervention_nodes}) and $G_{\doit(W)}$ (def \ref{def:G_hard_intervention}) given in the disambiguation paragraph above.

\smallskip\noindent\emph{Shorthand for the two single-step CDMGs.}
By def \ref{def:cdmg_intervention_nodes},
\[
    G_{\doit(I_{W_1})} \;=\; (J_1,\; V_1,\; E_1,\; L_1),
\]
with
\[
    J_1 := J \cup \lC I_w \st w \in W_1 \sm J \rC, \quad V_1 := V, \quad E_1 := E \cup \lC (I_w, w) \st w \in W_1 \sm J \rC, \quad L_1 := L;
\]
and by def \ref{def:G_hard_intervention},
\[
    G_{\doit(W_2)} \;=\; (J_2,\; V_2,\; E_2,\; L_2),
\]
with
\[
    J_2 := J \cup W_2, \quad V_2 := V \sm W_2, \quad E_2 := E \sm \lC (v_1, v_2) \in E \st v_2 \in W_2 \rC, \quad L_2 := L \sm \lC (v_1, v_2) \in L \st v_2 \in W_2 \rC,
\]
the bidirected-edge component $L_2$ read with the symmetrisation noted in the bullet on $G_{\doit(W)}$ in the statement block (so $L_2$ satisfies the symmetry axiom of def \ref{def-cdmg}).

\medskip\noindent\textbf{Proof of (a).}
We split sub-claim (a) into:
\begin{itemize}
    \item[(a-1)] $\lp G_{\doit(I_{W_1})} \rp_{\doit(I_{W_2})} \;=\; G_{\doit(I_{W_1 \cup W_2})}$;
    \item[(a-2)] $\lp G_{\doit(I_{W_2})} \rp_{\doit(I_{W_1})} \;=\; G_{\doit(I_{W_1 \cup W_2})}$.
\end{itemize}
We prove (a-1) by an explicit componentwise calculation followed by a carrier-map injectivity check; (a-2) is obtained from (a-1) by symmetry of the roles of $W_1$ and $W_2$ combined with commutativity of union; the LN's stated swap-symmetry $\lp G_{\doit(I_{W_1})} \rp_{\doit(I_{W_2})} = \lp G_{\doit(I_{W_2})} \rp_{\doit(I_{W_1})}$ then follows by transitivity of equality on the shared right-hand side $G_{\doit(I_{W_1 \cup W_2})}$.

\smallskip\noindent\emph{A preliminary carrier-set identity for (a-1): $W_2 \sm J_1 = W_2 \sm J$.}
The outer operation $\doit(I_{W_2})$ refers to the input CDMG $G_{\doit(I_{W_1})}$ through its $J$-component $J_1$, in the form $W_2 \sm J_1$ (the index set of fresh symbols and of fresh edges, per items i.\ and iii.\ of def \ref{def:cdmg_intervention_nodes}). Computing,
\[
    W_2 \cap J_1 \;=\; W_2 \cap \lp J \cup \lC I_w \st w \in W_1 \sm J \rC \rp \;=\; (W_2 \cap J) \cup \lp W_2 \cap \lC I_w \st w \in W_1 \sm J \rC \rp.
\]
By the freshness clause of def \ref{def:cdmg_intervention_nodes}, each fresh symbol $I_w$ (with $w \in W_1 \sm J$) is type-disjoint from $J \cup V$; since $W_2 \subseteq J \cup V$, the second intersection is empty, and $W_2 \cap J_1 = W_2 \cap J$. Hence $W_2 \sm J_1 = W_2 \sm J$. (Note that this identity does \emph{not} consume the disjointness hypothesis $W_1 \cap W_2 = \emptyset$ -- it follows from the freshness clause alone.)

\smallskip\noindent\emph{Componentwise verification of (a-1).}

\medskip\noindent\emph{Input nodes.}
By item i.\ of def \ref{def:cdmg_intervention_nodes} applied to the outer operation $\doit(I_{W_2})$ on the input CDMG $G_{\doit(I_{W_1})} = (J_1, V_1, E_1, L_1)$,
\[
    J_{\lp G_{\doit(I_{W_1})} \rp_{\doit(I_{W_2})}} \;=\; J_1 \cup \lC I_w \st w \in W_2 \sm J_1 \rC.
\]
Substituting $J_1 = J \cup \lC I_w \st w \in W_1 \sm J \rC$ and using $W_2 \sm J_1 = W_2 \sm J$ from the carrier identity above,
\[
    J_{\lp G_{\doit(I_{W_1})} \rp_{\doit(I_{W_2})}} \;=\; J \cup \lC I_w \st w \in W_1 \sm J \rC \cup \lC I_w \st w \in W_2 \sm J \rC.
\]
The disjointness hypothesis $W_1 \cap W_2 = \emptyset$ implies $(W_1 \sm J) \cap (W_2 \sm J) = \emptyset$, so the two index sets are disjoint; combined with the freshness clause of def \ref{def:cdmg_intervention_nodes} (each $w$ produces a distinct fresh symbol $I_w$), the two fresh-symbol sets $\lC I_w \st w \in W_1 \sm J \rC$ and $\lC I_w \st w \in W_2 \sm J \rC$ are themselves disjoint. Their union therefore reindexes as
\[
    \lC I_w \st w \in W_1 \sm J \rC \cup \lC I_w \st w \in W_2 \sm J \rC \;=\; \lC I_w \st w \in (W_1 \sm J) \cup (W_2 \sm J) \rC \;=\; \lC I_w \st w \in (W_1 \cup W_2) \sm J \rC,
\]
where the last equality is the set-theoretic identity $(A \sm C) \cup (B \sm C) = (A \cup B) \sm C$ instantiated at $A = W_1$, $B = W_2$, $C = J$. Substituting,
\[
    J_{\lp G_{\doit(I_{W_1})} \rp_{\doit(I_{W_2})}} \;=\; J \cup \lC I_w \st w \in (W_1 \cup W_2) \sm J \rC \;=\; J_{\doit(I_{W_1 \cup W_2})},
\]
the last equality being item i.\ of def \ref{def:cdmg_intervention_nodes} applied to the single-step operation $\doit(I_{W_1 \cup W_2})$ on $G$. \checkmark

\medskip\noindent\emph{Output nodes.}
By item ii.\ of def \ref{def:cdmg_intervention_nodes}, the operation $\doit(I_W)$ leaves $V$ unchanged for any $W \subseteq J \cup V$. Applied twice (first to $G$ with $W_1$, then to $G_{\doit(I_{W_1})}$ with $W_2$) and once on the right (to $G$ with $W_1 \cup W_2$),
\[
    V_{\lp G_{\doit(I_{W_1})} \rp_{\doit(I_{W_2})}} \;=\; V_1 \;=\; V \;=\; V_{\doit(I_{W_1 \cup W_2})}. \quad \checkmark
\]
(No use of disjointness is needed here -- $V$ is unchanged regardless of $W_1, W_2$.)

\medskip\noindent\emph{Directed edges.}
By item iii.\ of def \ref{def:cdmg_intervention_nodes} applied to the outer operation,
\[
    E_{\lp G_{\doit(I_{W_1})} \rp_{\doit(I_{W_2})}} \;=\; E_1 \cup \lC (I_w, w) \st w \in W_2 \sm J_1 \rC.
\]
Substituting $E_1 = E \cup \lC (I_w, w) \st w \in W_1 \sm J \rC$ and using $W_2 \sm J_1 = W_2 \sm J$,
\[
    E_{\lp G_{\doit(I_{W_1})} \rp_{\doit(I_{W_2})}} \;=\; E \cup \lC (I_w, w) \st w \in W_1 \sm J \rC \cup \lC (I_w, w) \st w \in W_2 \sm J \rC.
\]
By the disjointness hypothesis $W_1 \cap W_2 = \emptyset$ (yielding $(W_1 \sm J) \cap (W_2 \sm J) = \emptyset$) and the freshness clause of def \ref{def:cdmg_intervention_nodes} (each $w$ produces a distinct $I_w$, so distinct $w$'s give distinct ordered pairs $(I_w, w)$), the two fresh-edge sets are disjoint. Their union reindexes as
\[
    \lC (I_w, w) \st w \in W_1 \sm J \rC \cup \lC (I_w, w) \st w \in W_2 \sm J \rC \;=\; \lC (I_w, w) \st w \in (W_1 \cup W_2) \sm J \rC,
\]
by the same $(A \sm C) \cup (B \sm C) = (A \cup B) \sm C$ identity used in the input-node step. Substituting,
\[
    E_{\lp G_{\doit(I_{W_1})} \rp_{\doit(I_{W_2})}} \;=\; E \cup \lC (I_w, w) \st w \in (W_1 \cup W_2) \sm J \rC \;=\; E_{\doit(I_{W_1 \cup W_2})},
\]
the last equality being item iii.\ of def \ref{def:cdmg_intervention_nodes} applied to $\doit(I_{W_1 \cup W_2})$. \checkmark

\medskip\noindent\emph{Bidirected edges.}
By item iv.\ of def \ref{def:cdmg_intervention_nodes}, the operation $\doit(I_W)$ leaves $L$ unchanged for any $W \subseteq J \cup V$. Applied twice on the left and once on the right,
\[
    L_{\lp G_{\doit(I_{W_1})} \rp_{\doit(I_{W_2})}} \;=\; L_1 \;=\; L \;=\; L_{\doit(I_{W_1 \cup W_2})}. \quad \checkmark
\]
(No use of disjointness is needed here -- $L$ is unchanged regardless of $W_1, W_2$.)

\medskip\noindent\emph{Injectivity of the canonical flatten map $\flattenIntExt$ on the iterated extended graph's $J \cup V$.}
The canonical flatten map $\flattenIntExt$ between the iterated-extension carrier $J_{\lp G_{\doit(I_{W_1})} \rp_{\doit(I_{W_2})}} \cup V_{\lp G_{\doit(I_{W_1})} \rp_{\doit(I_{W_2})}}$ and the single-extension carrier $J_{\doit(I_{W_1 \cup W_2})} \cup V_{\doit(I_{W_1 \cup W_2})}$ sends every iterated-extension node to its single-step counterpart, via the four-cell piecewise rule
\begin{itemize}
    \item $\flattenIntExt(\iota(\iota(j))) \;:=\; \iota(j)$ for $j \in J$;
    \item $\flattenIntExt(\iota(\iota(v))) \;:=\; \iota(v)$ for $v \in V$;
    \item $\flattenIntExt(\iota(I_{w_1})) \;:=\; I_{w_1}$ for $w_1 \in W_1 \sm J$;
    \item $\flattenIntExt(I_{w_2}) \;:=\; I_{w_2}$ for $w_2 \in W_2 \sm J$,
\end{itemize}
where: $\iota$ denotes the unsplit-injection of def \ref{def:cdmg_intervention_nodes} carrying $J \cup V$ into the extended carrier (used twice on the iterated side and once on the single-step side); $\iota(I_{w_1})$ in the third case is the inner extension's fresh symbol $I_{w_1}$ (introduced by the inner $\doit(I_{W_1})$) carried up to the iterated carrier by the outer $\iota$ (per the LN's identification of carriers -- the outer $\iota$ acts as the identity on inner symbols not in $W_2$); and $I_{w_2}$ in the fourth case is the outer extension's fresh symbol introduced by $\doit(I_{W_2})$ on the inner CDMG $G_{\doit(I_{W_1})}$, indexed by $w_2 \in W_2 \sm J_1 = W_2 \sm J$ (using the carrier identity above to write the index range in terms of $J$ rather than $J_1$, and under the LN's identification $I_{\iota(w_2)} = I_{w_2}$ since $\iota(w_2) = w_2$ for $w_2 \in J \cup V$).

We verify that $\flattenIntExt$ is \emph{injective} on the iterated-extension's joint input-output set -- i.e., distinct elements of $J_{\lp G_{\doit(I_{W_1})} \rp_{\doit(I_{W_2})}} \cup V_{\lp G_{\doit(I_{W_1})} \rp_{\doit(I_{W_2})}}$ are sent to distinct elements of $J_{\doit(I_{W_1 \cup W_2})} \cup V_{\doit(I_{W_1 \cup W_2})}$ -- by a case analysis on the four-cell partition of the source.

Combining the input-node and output-node identities verified above with item i.\ of def \ref{def:cdmg_intervention_nodes}'s shape for the iterated $J_1$ and the carrier identity $W_2 \sm J_1 = W_2 \sm J$, the iterated graph's joint input-output set decomposes as the disjoint union of four cells
\begin{align*}
    J_{\lp G_{\doit(I_{W_1})} \rp_{\doit(I_{W_2})}} \cup V_{\lp G_{\doit(I_{W_1})} \rp_{\doit(I_{W_2})}}
        &\;=\; \iota(\iota(J)) \;\dcup\; \iota(\iota(V)) \\
        &\quad\dcup\; \iota\lp\lC I_{w_1} \st w_1 \in W_1 \sm J \rC\rp \\
        &\quad\dcup\; \lC I_{w_2} \st w_2 \in W_2 \sm J \rC,
\end{align*}
where:
\begin{itemize}
    \item cell (1) $\iota(\iota(J))$ is the doubly-unsplit copy of $J$, sitting inside the iterated $J$-component as the image of $J \hookrightarrow J_1 \hookrightarrow J_{\lp G_{\doit(I_{W_1})} \rp_{\doit(I_{W_2})}}$;
    \item cell (2) $\iota(\iota(V))$ is the doubly-unsplit copy of $V$, sitting inside the iterated $V$-component (which equals $V_1 = V$ verbatim, by item ii.\ of def \ref{def:cdmg_intervention_nodes} applied twice);
    \item cell (3) $\iota\lp\lC I_{w_1} \st w_1 \in W_1 \sm J \rC\rp$ is the inner extension's fresh symbols carried up to the iterated carrier by the outer $\iota$ (the inner $\doit(I_{W_1})$ produces these, and the outer $\doit(I_{W_2})$ leaves them in place via the outer $\iota$ since $I_{w_1} \notin W_2.\mathrm{image}(\iota) = W_2$ by the freshness clause);
    \item cell (4) $\lC I_{w_2} \st w_2 \in W_2 \sm J \rC$ is the outer extension's fresh symbols, indexed by $W_2 \sm J_1 = W_2 \sm J$ via the carrier identity above.
\end{itemize}

\emph{Pairwise disjointness of the four cells inside the iterated $J \cup V$.}
Cells (1) and (2) are disjoint by the disjointness axiom $J \cap V = \emptyset$ of def \ref{def-cdmg} together with $\iota \circ \iota$'s injectivity ($\iota$ being the identity-on-set inclusion of $J \cup V$ into the extended carrier, applied twice). Cells (1) and (2) are disjoint from cells (3) and (4) by the freshness clause of def \ref{def:cdmg_intervention_nodes}: the fresh symbols $\lC I_{w_1} \st w_1 \in W_1 \sm J \rC$ are type-disjoint from $J \cup V \supseteq \iota(\iota(J)) \cup \iota(\iota(V))$ (after applying the outer $\iota$ which leaves them as fresh symbols), and the outer fresh symbols $\lC I_{w_2} \st w_2 \in W_2 \sm J \rC$ are type-disjoint from the inner carrier $J_1 \cup V_1$ (the freshness clause applied to the outer operation), hence in particular from $\iota(\iota(J)) \cup \iota(\iota(V)) \subseteq \iota(J_1 \cup V_1)$. Finally, cells (3) and (4) are disjoint from each other because the inner fresh symbols and the outer fresh symbols are produced by two distinct applications of def \ref{def:cdmg_intervention_nodes}'s freshness clause -- the outer clause is applied to the inner CDMG $G_{\doit(I_{W_1})}$ as input and produces symbols type-disjoint from the inner CDMG's carrier $J_1 \cup V_1$, while the inner fresh symbols $\iota(I_{w_1})$ \emph{do} live in $\iota(J_1) \subseteq \iota(J_1 \cup V_1)$. Hence the four cells are pairwise disjoint inside the iterated joint input-output set.

\emph{Within-cell injectivity.}
On each of the four cells the map $\flattenIntExt$ acts by a single composed-constructor chain that is injective in the underlying $j$/$v$/$w_1$/$w_2$-label:
\begin{itemize}
    \item ``$\iota(\iota(j)) \mapsto \iota(j)$'' on cell (1): injective in $j$ because $\iota$ is the identity-on-set inclusion of $J \cup V$ into the extended carrier on both sides (def \ref{def:cdmg_intervention_nodes}, item i.), so the composition is the identity on $J$.
    \item ``$\iota(\iota(v)) \mapsto \iota(v)$'' on cell (2): same argument with $v$ in place of $j$, via item ii.\ of def \ref{def:cdmg_intervention_nodes}.
    \item ``$\iota(I_{w_1}) \mapsto I_{w_1}$'' on cell (3): injective in $w_1$ because the fresh symbols $\lC I_{w_1} \st w_1 \in W_1 \sm J \rC$ are pairwise distinct (def \ref{def:cdmg_intervention_nodes}'s freshness clause: distinct $w_1$'s produce distinct $I_{w_1}$'s), and $\iota$ is injective.
    \item ``$I_{w_2} \mapsto I_{w_2}$'' on cell (4): injective in $w_2$ because the outer fresh symbols are pairwise distinct in their index $w_2$ (same freshness clause, applied to the outer operation), and the single-step fresh symbols $I_{w_2}$ on the codomain side are also pairwise distinct in $w_2$.
\end{itemize}

\emph{Across-cell injectivity.}
It remains to show that no collision occurs between two elements drawn from two different cells of the four-cell partition. There are $\binom{4}{2} = 6$ structural cross-cell collision patterns to rule out; we walk through each.
\begin{itemize}
    \item \emph{Cell (1) vs cell (2).} Images $\iota(j)$ vs $\iota(v)$ for some $j \in J$, $v \in V$. An equality $\iota(j) = \iota(v)$ in the single-step carrier forces $j = v$ by injectivity of $\iota$ (its identity-on-set realization on $J \cup V$), hence $j = v \in J \cap V = \emptyset$ (def \ref{def-cdmg}). No collision.
    \item \emph{Cell (1) vs cell (3).} Images $\iota(j)$ vs $I_{w_1}$ for some $j \in J$, $w_1 \in W_1 \sm J$. By the freshness clause of def \ref{def:cdmg_intervention_nodes} applied to $\doit(I_{W_1 \cup W_2})$ on $G$, the fresh symbol $I_{w_1} \in \lC I_w \st w \in (W_1 \cup W_2) \sm J \rC$ is type-disjoint from $J \cup V \supseteq \iota(J)$. No collision.
    \item \emph{Cell (1) vs cell (4).} Images $\iota(j)$ vs $I_{w_2}$ for some $j \in J$, $w_2 \in W_2 \sm J$. Same argument as cells (1) vs (3), with $w_2$ in place of $w_1$. No collision.
    \item \emph{Cell (2) vs cell (3).} Images $\iota(v)$ vs $I_{w_1}$ for some $v \in V$, $w_1 \in W_1 \sm J$. Same freshness argument, with $v \in V \subseteq J \cup V$. No collision.
    \item \emph{Cell (2) vs cell (4).} Images $\iota(v)$ vs $I_{w_2}$ for some $v \in V$, $w_2 \in W_2 \sm J$. Same. No collision.
    \item \emph{Cell (3) vs cell (4).} Images $I_{w_1}$ vs $I_{w_2}$ for some $w_1 \in W_1 \sm J$, $w_2 \in W_2 \sm J$. An equality $I_{w_1} = I_{w_2}$ in the single-step carrier forces $w_1 = w_2$ by injectivity of the assignment $w \mapsto I_w$ on the single-step source set $(W_1 \cup W_2) \sm J$ (the freshness clause of def \ref{def:cdmg_intervention_nodes} produces a distinct fresh symbol for each distinct $w$). But $w_1 \in W_1 \sm J \subseteq W_1$ and $w_2 \in W_2 \sm J \subseteq W_2$, and the disjointness hypothesis $W_1 \cap W_2 = \emptyset$ rules out $w_1 = w_2$. No collision. \emph{This is the load-bearing use of the disjointness hypothesis $W_1 \cap W_2 = \emptyset$.}
\end{itemize}
All six cross-cell collision patterns have been ruled out, so the canonical flatten map $\flattenIntExt$ is injective on $J_{\lp G_{\doit(I_{W_1})} \rp_{\doit(I_{W_2})}} \cup V_{\lp G_{\doit(I_{W_1})} \rp_{\doit(I_{W_2})}}$. \checkmark

\smallskip
Combining the four componentwise equalities above with the carrier-map injectivity check just established, the CDMGs $\lp G_{\doit(I_{W_1})} \rp_{\doit(I_{W_2})}$ and $G_{\doit(I_{W_1 \cup W_2})}$ agree field by field on each of $(J, V, E, L)$ -- with the four image-equality identities understood modulo the canonical flatten map $\flattenIntExt$, which is in turn injective on the iterated graph's joint input-output set -- and are therefore equal as CDMGs (def \ref{def-cdmg}) in the strengthened componentwise-up-to-carrier-bijection reading recorded in the statement block above. This establishes (a-1).

\medskip\noindent\emph{Proof of (a-2) by symmetry of the roles of $W_1$ and $W_2$.}
The five-component verification of (a-1) above (four image-equality identities plus the carrier-map injectivity check) depends on the pair $(W_1, W_2)$ only through their symbolic roles in: the carrier identity $W_2 \sm J_1 = W_2 \sm J$; the disjointness hypothesis $W_1 \cap W_2 = \emptyset$ (consumed in the input-node, directed-edge, and cell-(3)-vs-cell-(4) of the injectivity check); and the freshness-clause arguments (whose roles in $W_1$ and $W_2$ are interchangeable). Each is preserved under the swap $W_1 \leftrightarrow W_2$: the carrier identity becomes $W_1 \sm J_1' = W_1 \sm J$ where $J_1' := J \cup \lC I_w \st w \in W_2 \sm J \rC$ is the $J$-component of $G_{\doit(I_{W_2})}$ (the same argument applies, using the freshness clause of def \ref{def:cdmg_intervention_nodes}); the disjointness hypothesis is $W_2 \cap W_1 = \emptyset$, identical to $W_1 \cap W_2 = \emptyset$ by symmetry of $\cap$.

\medskip\noindent\emph{Injectivity of the canonical flatten map $\flattenIntExt$ on the iterated extended graph's $J \cup V$ (symmetric counterpart).}
For sub-claim (a-2)'s iterated graph $\lp G_{\doit(I_{W_2})} \rp_{\doit(I_{W_1})}$, the carrier-map injectivity check is the same as in (a-1) with the substitution $W_1 \leftrightarrow W_2$ throughout: the iterated graph's joint input-output set decomposes as the four-cell disjoint union
\begin{align*}
    J_{\lp G_{\doit(I_{W_2})} \rp_{\doit(I_{W_1})}} \cup V_{\lp G_{\doit(I_{W_2})} \rp_{\doit(I_{W_1})}}
        &\;=\; \iota(\iota(J)) \;\dcup\; \iota(\iota(V)) \\
        &\quad\dcup\; \iota\lp\lC I_{w_2} \st w_2 \in W_2 \sm J \rC\rp \\
        &\quad\dcup\; \lC I_{w_1} \st w_1 \in W_1 \sm J \rC,
\end{align*}
with the inner-and-outer roles of $W_1$ and $W_2$ exchanged (cells (3) and (4) are now indexed by $W_2 \sm J$ and $W_1 \sm J$ respectively, since the inner extension is $\doit(I_{W_2})$ and the outer is $\doit(I_{W_1})$). The pairwise-disjointness argument for the four cells is the same as in (a-1) -- using the freshness clause of def \ref{def:cdmg_intervention_nodes} (twice) and the $J \cap V = \emptyset$ axiom of def \ref{def-cdmg}. The within-cell injectivity argument is unchanged (each of the four cells maps injectively in its underlying $j$/$v$/$w_2$/$w_1$-label). Across-cell injectivity: the cell-(1)-vs-(2) collision is ruled out by $J \cap V = \emptyset$; the cell-(1)/(2)-vs-(3)/(4) collisions are ruled out by the freshness clause; and the cell-(3)-vs-(4) collision $I_{w_2} = I_{w_1}$ with $w_2 \in W_2 \sm J$ and $w_1 \in W_1 \sm J$ forces $w_2 = w_1$ by injectivity of $w \mapsto I_w$, ruled out by the disjointness hypothesis $W_2 \cap W_1 = W_1 \cap W_2 = \emptyset$ (symmetric in the pair, identical to the (a-1) case). Hence $\flattenIntExt$ is injective on $J_{\lp G_{\doit(I_{W_2})} \rp_{\doit(I_{W_1})}} \cup V_{\lp G_{\doit(I_{W_2})} \rp_{\doit(I_{W_1})}}$. \checkmark

\medskip
Re-running the five-component verification (four image-equality identities plus the carrier-map injectivity check) with the substitution $W_1 \leftrightarrow W_2$ throughout therefore yields
\[
    \lp G_{\doit(I_{W_2})} \rp_{\doit(I_{W_1})} \;=\; G_{\doit(I_{W_2 \cup W_1})}.
\]
The right-hand sides $G_{\doit(I_{W_2 \cup W_1})}$ and $G_{\doit(I_{W_1 \cup W_2})}$ coincide as CDMGs: def \ref{def:cdmg_intervention_nodes} depends on its argument $W$ only through the underlying set $W \subseteq J \cup V$ (each of items i.--iv.\ refers to $W$ only via the operations $J \cup \cdot$, the set difference $\cdot \sm J$, and the membership predicate $\cdot \in W$), so equal underlying sets produce equal CDMGs componentwise; and $W_2 \cup W_1 = W_1 \cup W_2$ as sets by commutativity of union. Chaining,
\[
    \lp G_{\doit(I_{W_2})} \rp_{\doit(I_{W_1})} \;=\; G_{\doit(I_{W_2 \cup W_1})} \;=\; G_{\doit(I_{W_1 \cup W_2})},
\]
which establishes (a-2).

\medskip\noindent\emph{Swap-symmetry and recovery of the LN's triple equality.}
By (a-1) and (a-2), both iterated operations equal the single CDMG $G_{\doit(I_{W_1 \cup W_2})}$; transitivity of equality through this shared right-hand side then yields the LN's stated triple equality
\[
    \lp G_{\doit(I_{W_1})} \rp_{\doit(I_{W_2})} \;=\; G_{\doit(I_{W_1 \cup W_2})} \;=\; \lp G_{\doit(I_{W_2})} \rp_{\doit(I_{W_1})},
\]
and in particular the swap-symmetry reading $\lp G_{\doit(I_{W_1})} \rp_{\doit(I_{W_2})} = \lp G_{\doit(I_{W_2})} \rp_{\doit(I_{W_1})}$. This establishes sub-claim (a).

\medskip\noindent\textbf{Proof of (b).}
We compute each of the two iterated CDMGs $\lp G_{\doit(I_{W_1})} \rp_{\doit(W_2)}$ (henceforth ``LHS'') and $\lp G_{\doit(W_2)} \rp_{\doit(I_{W_1})}$ (henceforth ``middle'') componentwise into the same four-tuple, equate them by componentwise agreement, and then identify their common value with $G_{\doit(I_{W_1}, W_2)}$ by appeal to the defining equation of the mixed-argument notation,
\[
    G_{\doit(I_{W_1}, W_2)} \;:=\; \lp G_{\doit(I_{W_1})} \rp_{\doit(W_2)},
\]
introduced in the statement block above. Throughout, we keep the shorthand $(J_1, V_1, E_1, L_1)$ for $G_{\doit(I_{W_1})}$ and $(J_2, V_2, E_2, L_2)$ for $G_{\doit(W_2)}$ fixed at the start of the proof. (Sub-claim (b) has no carrier mismatch -- both LHS and middle live on the same single-extension carrier $J_{\doit(I_{W_1})} \cup V_{\doit(I_{W_1})}$ -- so no carrier-bijection identification is needed and the four componentwise equalities are read literally, per the statement-block carrier-mismatch paragraph.)

\smallskip\noindent\emph{A carrier-set identity for the middle ordering: $W_1 \sm J_2 = W_1 \sm J$.}
Since $J_2 = J \cup W_2$,
\[
    W_1 \sm J_2 \;=\; W_1 \sm (J \cup W_2) \;=\; (W_1 \sm J) \sm W_2.
\]
By the disjointness hypothesis $W_1 \cap W_2 = \emptyset$, $(W_1 \sm J) \cap W_2 \subseteq W_1 \cap W_2 = \emptyset$, so the removal of $W_2$ from $W_1 \sm J$ has no effect; hence $(W_1 \sm J) \sm W_2 = W_1 \sm J$ and therefore $W_1 \sm J_2 = W_1 \sm J$. This identity is used in the middle-ordering input-node and directed-edge calculations below; it is the (only) place where disjointness is consumed in those calculations.

\smallskip\noindent\emph{Componentwise calculation of the LHS, $\lp G_{\doit(I_{W_1})} \rp_{\doit(W_2)}$.}
Applying $\doit(W_2)$ to $G_{\doit(I_{W_1})} = (J_1, V_1, E_1, L_1)$ per def \ref{def:G_hard_intervention} (well-typedness $W_2 \subseteq J \cup V \subseteq J_1 \cup V_1$ from the statement block):

\medskip\noindent\emph{Input nodes (item i.\ of def \ref{def:G_hard_intervention}):}
\begin{align*}
    J_{\mathrm{LHS}}
        &\;=\; J_1 \cup W_2 \\
        &\;=\; \lp J \cup \lC I_w \st w \in W_1 \sm J \rC \rp \cup W_2 \\
        &\;=\; (J \cup W_2) \cup \lC I_w \st w \in W_1 \sm J \rC,
\end{align*}
by associativity and commutativity of $\cup$. (No use of disjointness here.)

\medskip\noindent\emph{Output nodes (item ii.\ of def \ref{def:G_hard_intervention}):}
\[
    V_{\mathrm{LHS}} \;=\; V_1 \sm W_2 \;=\; V \sm W_2.
\]
(No use of disjointness here.)

\medskip\noindent\emph{Directed edges (item iii.\ of def \ref{def:G_hard_intervention}):}
\[
    E_{\mathrm{LHS}} \;=\; E_1 \sm \lC (v_1, v_2) \in E_1 \st v_2 \in W_2 \rC.
\]
Substituting $E_1 = E \cup \lC (I_w, w) \st w \in W_1 \sm J \rC$ and splitting the removal over the union (using the elementary identity $(A \cup B) \sm S = (A \sm S) \cup (B \sm S)$ applied with $S$ the removal predicate's set), we get
\[
    E_{\mathrm{LHS}} \;=\; \lp E \sm \lC (v_1, v_2) \in E \st v_2 \in W_2 \rC \rp \cup \lp \lC (I_w, w) \st w \in W_1 \sm J \rC \sm \lC (v_1, v_2) \st v_2 \in W_2 \rC \rp.
\]
For the second piece: every fresh edge $(I_w, w)$ in $\lC (I_w, w) \st w \in W_1 \sm J \rC$ has its head $v_2 = w \in W_1 \sm J \subseteq W_1$; by the disjointness hypothesis $W_1 \cap W_2 = \emptyset$, $w \notin W_2$, so the predicate ``$v_2 \in W_2$'' is never satisfied on the fresh edges. Hence the second piece collapses to $\lC (I_w, w) \st w \in W_1 \sm J \rC$ unchanged, and
\[
    E_{\mathrm{LHS}} \;=\; \lp E \sm \lC (v_1, v_2) \in E \st v_2 \in W_2 \rC \rp \cup \lC (I_w, w) \st w \in W_1 \sm J \rC.
\]
This is the (only) place where disjointness is consumed in the LHS calculation.

\medskip\noindent\emph{Bidirected edges (item iv.\ of def \ref{def:G_hard_intervention}, with the symmetrisation noted in the statement block):}
\[
    L_{\mathrm{LHS}} \;=\; L_1 \sm \lC (v_1, v_2) \in L_1 \st v_2 \in W_2 \rC \;=\; L \sm \lC (v_1, v_2) \in L \st v_2 \in W_2 \rC,
\]
since $L_1 = L$. The symmetrisation noted in the statement block applies here (and only here, on the bidirected side): $L_{\mathrm{LHS}}$ is taken to be the symmetric closure of the displayed one-sided removal, which is what the CDMG axioms (def \ref{def-cdmg}) require. (No use of disjointness here.)

\smallskip\noindent\emph{Componentwise calculation of the middle, $\lp G_{\doit(W_2)} \rp_{\doit(I_{W_1})}$.}
Applying $\doit(I_{W_1})$ to $G_{\doit(W_2)} = (J_2, V_2, E_2, L_2)$ per def \ref{def:cdmg_intervention_nodes} (well-typedness $W_1 \subseteq J \cup V = J_2 \cup V_2$ from the statement block):

\medskip\noindent\emph{Input nodes (item i.\ of def \ref{def:cdmg_intervention_nodes}):}
\[
    J_{\mathrm{mid}} \;=\; J_2 \cup \lC I_w \st w \in W_1 \sm J_2 \rC.
\]
Substituting $J_2 = J \cup W_2$ and using the carrier identity $W_1 \sm J_2 = W_1 \sm J$ (which consumes disjointness),
\[
    J_{\mathrm{mid}} \;=\; (J \cup W_2) \cup \lC I_w \st w \in W_1 \sm J \rC \;=\; J_{\mathrm{LHS}}. \quad \checkmark
\]

\medskip\noindent\emph{Output nodes (item ii.\ of def \ref{def:cdmg_intervention_nodes}):}
\[
    V_{\mathrm{mid}} \;=\; V_2 \;=\; V \sm W_2 \;=\; V_{\mathrm{LHS}}. \quad \checkmark
\]
(The operation $\doit(I_{W_1})$ leaves $V_2$ unchanged; no use of disjointness here.)

\medskip\noindent\emph{Directed edges (item iii.\ of def \ref{def:cdmg_intervention_nodes}):}
\[
    E_{\mathrm{mid}} \;=\; E_2 \cup \lC (I_w, w) \st w \in W_1 \sm J_2 \rC.
\]
Substituting $E_2 = E \sm \lC (v_1, v_2) \in E \st v_2 \in W_2 \rC$ and using the carrier identity $W_1 \sm J_2 = W_1 \sm J$ (which consumes disjointness),
\[
    E_{\mathrm{mid}} \;=\; \lp E \sm \lC (v_1, v_2) \in E \st v_2 \in W_2 \rC \rp \cup \lC (I_w, w) \st w \in W_1 \sm J \rC \;=\; E_{\mathrm{LHS}}. \quad \checkmark
\]

\medskip\noindent\emph{Bidirected edges (item iv.\ of def \ref{def:cdmg_intervention_nodes}):}
\[
    L_{\mathrm{mid}} \;=\; L_2 \;=\; L \sm \lC (v_1, v_2) \in L \st v_2 \in W_2 \rC \;=\; L_{\mathrm{LHS}}. \quad \checkmark
\]
The outer operation $\doit(I_{W_1})$ leaves $L_2$ untouched (item iv.\ of def \ref{def:cdmg_intervention_nodes}), so the bidirected-edge content of the middle is entirely determined by the inner hard intervention $\doit(W_2)$. The LHS bidirected content was likewise entirely determined by the outer hard intervention $\doit(W_2)$ applied to $L_1 = L$, which produces the same set (with the same symmetrised reading carried by the LN's $L$-removal). Hence $L_{\mathrm{mid}} = L_{\mathrm{LHS}}$. (No use of disjointness here: the bidirected-edge component is unaffected by the doit$(I_{W_1})$ side and only sees the doit$(W_2)$ side, which removes the same set of bidirected edges incident to $W_2$ in both orderings.)

\smallskip\noindent\emph{Equating LHS and middle.}
By the four componentwise checks above, $J_{\mathrm{LHS}} = J_{\mathrm{mid}}$, $V_{\mathrm{LHS}} = V_{\mathrm{mid}}$, $E_{\mathrm{LHS}} = E_{\mathrm{mid}}$, and $L_{\mathrm{LHS}} = L_{\mathrm{mid}}$. Hence the two iterated CDMGs are equal as CDMGs (def \ref{def-cdmg}):
\[
    \lp G_{\doit(I_{W_1})} \rp_{\doit(W_2)} \;=\; \lp G_{\doit(W_2)} \rp_{\doit(I_{W_1})}.
\]

\smallskip\noindent\emph{Identification with the mixed-notation CDMG $G_{\doit(I_{W_1}, W_2)}$.}
The mixed-notation CDMG $G_{\doit(I_{W_1}, W_2)}$ was introduced in the statement block above by the defining equation
\[
    G_{\doit(I_{W_1}, W_2)} \;:=\; \lp G_{\doit(I_{W_1})} \rp_{\doit(W_2)},
\]
read in this exact ordering (intervention-node addition first, hard intervention second). Hence
\[
    \lp G_{\doit(I_{W_1})} \rp_{\doit(W_2)} \;=\; G_{\doit(I_{W_1}, W_2)}
\]
holds \emph{by definition} of the mixed-notation CDMG. Combining with the LHS\,$=$\,middle equality just established and transitivity,
\[
    \lp G_{\doit(W_2)} \rp_{\doit(I_{W_1})} \;=\; \lp G_{\doit(I_{W_1})} \rp_{\doit(W_2)} \;=\; G_{\doit(I_{W_1}, W_2)}.
\]
This is the LN's stated triple equality
\[
    \lp G_{\doit(I_{W_1})} \rp_{\doit(W_2)} \;=\; \lp G_{\doit(W_2)} \rp_{\doit(I_{W_1})} \;=\; G_{\doit(I_{W_1}, W_2)},
\]
and establishes sub-claim (b).
\end{proof}

\end{document}
